\section{Introducción}
\label{sec:intro}

El proyecto parte de un dataset sintético del sector salud y bienestar
distribuido como punto de partida para la entrega final. Reproduce un
escenario retail clásico ($50$ productos, una red de tiendas, catálogo de
promociones, devoluciones con motivo) sobre una ventana de $72$ meses
($2020$--$2025$, ambos inclusive). Los datos crudos están cargados en el
esquema \texttt{public} de PostgreSQL~16 en Docker (\texttt{localhost:5433},
BD \texttt{saleshealth}). La tabla~\ref{tab:cifras-clave} reúne las cifras
verificadas contra el DWH y los \emph{parquet} de analítica.

\begin{table}[H]
\centering
\caption{Cifras clave verificadas (DWH + \texttt{data/processed/}).}
\label{tab:cifras-clave}
\small
\begin{tabular}{lr@{\hspace{1.4em}}lr}
\toprule
Productos          & $50$            & Líneas de venta   & $42\,555$       \\
Clientes           & $5\,750$        & Líneas devueltas  & $2\,330$        \\
Tickets de venta   & $20\,000$       & Ingresos brutos   & $9{,}68$\,M\texteuro \\
Ingresos netos     & $9{,}40$\,M\texteuro & Ventana          & $72$ meses    \\
\bottomrule
\end{tabular}
\end{table}

\noindent\textbf{Objetivos.} (i) Diseñar un modelo Entidad--Relación sobre
\texttt{public} y un modelo dimensional en estrella sobre \texttt{dwh}.
(ii) Implementar un ETL reproducible \texttt{public$\to$staging$\to$dwh}.
(iii) Calcular el CLTV por cliente,
$\mathrm{CLTV}=\text{Ingresos}_t \times \text{Margen}_t \times
\text{Frecuencia}_t \times R_t$, junto con \emph{Recencia}, \emph{AOV} y
\emph{tasa de devolución}. (iv) Segmentar la base mediante \emph{PCA} +
\emph{KMeans}. (v) Documentar todo en un informe y un dashboard.

\noindent\textbf{Stack.} Python~3.11 en entorno \texttt{conda} \texttt{UAX};
PostgreSQL~16; \texttt{pandas}, \texttt{SQLAlchemy}~2.0, \texttt{scikit-learn};
notebooks Jupyter como orquestador; Flask para el dashboard; \texttt{pytest}.
